\chapter{符号表与主要定理依赖表：习题解答}

\begin{solution}{E.1}
\(A=[0,1]\) 时 \(\sup A=\max A=1\)。\(A=(0,1)\) 时只有上确界 \(1\)，没有最大元。\(A=\R\) 在实数范围内无上界，因而没有实数上确界，也没有最大元。下确界对应例可取 \((0,1]\)，其下确界为 \(0\) 而无最小元。
\end{solution}

\begin{solution}{E.2}
数列极限为
\[
 \forall\varepsilon>0\ \exists N\ \forall n\ge N.
\]
函数删心极限为
\[
 \forall\varepsilon>0\ \exists\delta>0\ \forall x\in D,\quad
 0<\abs{x-a}<\delta\Rightarrow\cdots.
\]
点连续把目标改为 \(f(a)\) 并允许 \(x=a\)。一致连续为
\[
 \forall\varepsilon>0\ \exists\delta>0\ \forall x,y\in D.
\]
逐点连续的 \(\delta\) 可依赖考察点，一致连续的 \(\delta\) 不能依赖 \(x,y\)。
\end{solution}

\begin{solution}{E.3}
对任意映射 \(f:X\to Y\) 和集合 \(A\subset Y\)，
\[
 f^{-1}(A)=\{x\in X:f(x)\in A\}
\]
始终是原像，不要求 \(f\) 可逆。只有当 \(f\) 为双射时，\(f^{-1}:Y\to X\) 才表示逆函数。可写 \(f^{[-1]}(A)\) 或明确写“集合原像”来避免同一段落中的混淆。
\end{solution}

\begin{solution}{E.4}
一条无循环路线是
\[
 \text{确界原理}\Rightarrow
 \text{单调有界收敛}\Rightarrow
 \text{区间套}\Rightarrow
 \text{Bolzano--Weierstrass}\Rightarrow
 \text{数列 Cauchy 完备}.
\]
函数 Cauchy 准则再选一条逼近考察点的数列，把局部 Cauchy 条件转成像 Cauchy 列；实数完备性产生候选极限，并由局部条件推广到整个删心邻域。
\end{solution}

\begin{solution}{E.5}
Heine 定理：函数极限等价于保持所有逼近点列的极限。Heine--Borel：实线中闭且有界等价于开覆盖有限化。Heine--Cantor：紧集上的连续函数一致连续。三者分别属于极限序列刻画、紧致集合刻画和连续性的全局化，名称相近但对象与结论不同。
\end{solution}

\begin{solution}{E.6}
在 \(\R\) 中有
\[
 \text{闭且有界}\Longleftrightarrow
 \text{开覆盖紧致}\Longleftrightarrow
 \text{序列紧致}.
\]
证明使用实数的有序与完备结构。一般度量空间中闭且有界未必紧致；一般拓扑空间中序列也未必刻画闭包或紧致性。因此后文只能保留以开覆盖定义的紧致性，其他等价需附空间条件。
\end{solution}

\begin{solution}{E.7}
速查表可按三列组织：输入对象如 \(f:D\to\R\)、数列或紧集；关键假设如聚点、完备值域、连续、单调、紧致；输出如极限存在、左右极限、介值、极值或一致连续。\chapterref{chap:function-limits}以极限类型为主，\chapterref{chap:continuity-discontinuity}以点连续与单调结构为主，\chapterref{chap:continuous-functions-closed-interval}以紧致和连通为主，\chapterref{chap:uniform-continuity}以统一误差与延拓为主。箭头必须指向正文中已经证明的定理。
\end{solution}

\begin{solution}{E.8}
以测度与概率记号为例，定稿前至少要统一：底层可测空间三元组的字体；Borel \(\sigma\)-代数记号；几乎处处与几乎必然的缩写；期望与积分的并用规则；依测度、\(L^p\)、几乎处处和分布收敛的箭头标记。这些属于全书级一致性决策，应提交宏观设计对话而不在具体内容文件中先行冻结。
\end{solution}
